\section*{Chapitre \ref{ch:m3:perpetual-futures} --- Contrats perpétuels}

\begin{solution}{exo:m3:perpetual-futures:1}
Le taux est positif, donc les longs paient les shorts $2 \times 60\,000 \times 0.03\% = \$36$.
\end{solution}

\begin{solution}{exo:m3:perpetual-futures:2}
0.02~\% est dans la zone morte ($-0.04\%$ à $+0.06\%$)~: financement de 0.01~\%. Pour 0.10~\%, $0.10\% - 0.05\% = 0.05\%$.
Pour $-0.08\%$, $-0.08\% + 0.05\% = -0.03\%$.
\end{solution}

\begin{solution}{exo:m3:perpetual-futures:3}
Une seule transaction, éventuellement petite ou manipulée, pourrait faire bouger le dernier cours et déclencher des liquidations~; la médiane
de l'estimation du financement fondée sur l'indice, de l'estimation fondée sur l'indice plus la base moyenne et de la dernière transaction ignore une valeur aberrante.
\end{solution}

\begin{solution}{exo:m3:perpetual-futures:4}
À 50\,000~: $60\,000\,(1/60\,000 - 1/50\,000) = -0.2$ bitcoin, $-\$10\,000$. À 75\,000~: $+0.2$ bitcoin,
$+\$15\,000$. En dollars, le P\&L est linéaire, $60\,000\,(S/60\,000 - 1)$~; en bitcoin, il ne l'est pas.
\end{solution}

\begin{solution}{exo:m3:perpetual-futures:5}
$10\,000 \times 0.000001 \times 300 = 3$ bitcoin, 180\,000~\$ à 60\,000~; la valeur en dollars serait différente à un autre
cours du bitcoin, alors que le mouvement de l'ether est le même.
\end{solution}

\begin{solution}{exo:m3:perpetual-futures:6}
$0.01\% \times 3 \times 365 = 10.95\%$~: environ 1.095~million de dollars.
\end{solution}

\begin{solution}{exo:m3:perpetual-futures:7}
12.93~\% (2020), 23.20~\% (2021), 2.95~\% (2022), 5.79~\% (2023), 8.89~\% (2024), 3.70~\% (2025) et 2.02~\% (2026,
au taux moyen jusqu'en septembre). Inférieur à 5~\% en 2022, 2025 et 2026.
\end{solution}

\begin{solution}{exo:m3:perpetual-futures:8}
Les variations de prix se compensent, mais le financement peut devenir négatif, la marge de la jambe short peut s'épuiser lors d'une hausse et être liquidée,
les plateformes qui détiennent les deux jambes peuvent faire défaut, la marge en \omterm{def:m3:blockchains-for-traders:stable}{stablecoin} peut subir un \omterm{def:m3:blockchains-for-traders:stable}{décrochage} (depeg), et la base à la sortie peut différer de celle à l'entrée.
\end{solution}

\begin{solution}{pb:m3:perpetual-futures:1}
\textbf{1.} 11.92~millions de dollars. \textbf{2.} $2 \times (0.10\% + 0.05\%) \times 100$ millions~: 0.30~million de dollars.
\textbf{3.} 10.95~\% par an. \textbf{4.} La prime moyenne a passé la plupart des intervalles dans la zone morte, où le financement
égale la composante d'intérêt. \textbf{5.} 4.16~millions de dollars. \textbf{6.} $30\% + 0.5\% \times 1.3 = 30.65\%$ du
notionnel~: 30.65~millions de dollars. \textbf{7.} 130.65~millions de dollars~; $11.62/130.65 = 8.89\%$. \textbf{8.} 11.62~\%.
\textbf{9.} Le short est liquidé lorsque sa perte ne laisse que la marge de maintenance, pour une hausse de $(30.65 -
0.5)/100.5 = 30.0\%$~; le fonds reste long en spot et non couvert, n'ayant rien vendu à prix élevé. \textbf{10.} Le spot est placé
dans un compte distinct~; son gain est latent et n'est pas utilisable comme marge, sauf si la plateforme le prend en compte.
\textbf{11.} Le short paie le financement tant que la position est couverte~; le revenu devient un coût. \textbf{12.} Les deux jambes et la marge
sont réparties sur des plateformes~: leur défaillance entraîne une perte pouvant aller jusqu'à la totalité du capital. \textbf{13.} La marge et le P\&L du short sont en
USDT~; un \omterm{def:m3:blockchains-for-traders:stable}{décrochage} modifie leur valeur en dollars. \textbf{14.} Si le \omterm{def:m3:perpetual-futures:perp}{contrat perpétuel} se traite moins cher que le spot à la sortie, clôturer le
short coûte plus que ce que rapporte la vente du spot par rapport à l'entrée~; la prime à l'entrée et à la sortie s'ajoute au résultat.
\textbf{15.} Le revenu de la stratégie est le plus élevé juste avant les krachs, lorsque les marges explosent, que les liquidations s'enchaînent et que le financement
peut s'inverser~: ses rendements et ses risques arrivent ensemble. \textbf{16.} Pas strictement~: elle rapporte un flux incertain et
supporte des risques de marge, de plateforme et de \omterm{def:m3:blockchains-for-traders:stable}{stablecoin}~; c'est un carry trade avec couverture. \textbf{17.} Les longs paient pour une
exposition à effet de levier sans emprunter ni détenir le spot~; le financement est le prix de cet effet de levier. \textbf{18.} Au moyen de limites de crédit par plateforme,
selon la profondeur et le financement de chaque plateforme, et afin que les appels de marge lors d'une hausse puissent être satisfaits sur toutes les plateformes à la fois.
\textbf{19.} \emph{Résultat à retenir~:} 8.89~\% sur le capital en 2024, avec 30.65~\% du notionnel en marge pour survivre à une
hausse de 30~\%. \textbf{20.} Le droit de détenir une position longue à effet de levier sans échéance, pour huit heures supplémentaires.
\end{solution}

\begin{solution}{iq:m3:perpetual-futures:1}
Un contrat à terme sans échéance~; les détenteurs échangent des paiements de financement à chaque intervalle, les longs versant aux shorts lorsque le contrat est
au-dessus de l'indice, et l'inverse lorsqu'il est en dessous, ce qui récompense les arbitragistes qui le ramènent vers l'indice.
\iqlookfor{le financement comme mécanisme de convergence.}
\end{solution}

\begin{solution}{iq:m3:perpetual-futures:2}
La dernière transaction est facile à faire bouger un instant~; liquider sur ce cours récompenserait la manipulation. Le \omterm{def:m3:perpetual-futures:index}{prix de marquage} (mark price) est construit à partir
de l'indice (plafonné, multiplateforme) et d'une base lissée, par exemple la médiane d'un prix indice plus base de financement,
d'un prix indice plus base à moyenne mobile et de la dernière transaction.
\iqlookfor{résistance à la manipulation et médiane.}
\end{solution}

\begin{solution}{iq:m3:perpetual-futures:3}
Linéaire~: P\&L dans la devise de cotation, option par défaut des traders raisonnant en dollars. Inverse~: notionnel en dollars, marge en coin et
règlement-livraison~; le P\&L en coin est concave, et un short adossé à une marge en coin constitue une position stable en dollars, utile aux
détenteurs de coin, comme les mineurs. Quanto~: exposition à un actif, réglée en un autre actif à taux fixe, ce qui ajoute l'actif de règlement
et la corrélation au risque.
\iqlookfor{devise du P\&L et qui détient naturellement chacun.}
\end{solution}

\begin{solution}{iq:m3:perpetual-futures:4}
Parce que la prime moyenne se situe dans la zone morte, où le clamp fixe le financement à la composante d'intérêt de
0.01~\%, quelle que soit la prime, tant qu'elle reste à moins de 0.05~\% de celle-ci.
\iqlookfor{le clamp.}
\end{solution}

\begin{solution}{iq:m3:perpetual-futures:5}
Financement négatif~; appels de marge sur les jambes short lors d'une hausse, éventuellement sur plusieurs plateformes à la fois, et liquidation si
ceux-ci restent impayés~; défaillance de plateforme~; \omterm{def:m3:blockchains-for-traders:stable}{décrochage} du \omterm{def:m3:blockchains-for-traders:stable}{stablecoin}~; variations de la base au dénouement~; risque opérationnel lors du transfert de collatéral. Limites par
plateforme, coussins de marge, marge croisée lorsque c'est possible, et un plan de dénouement.
\iqlookfor{la liquidité de la marge comme principal risque.}
\end{solution}

\begin{solution}{iq:m3:perpetual-futures:6}
Deux heures avant, les trois quarts des échantillons sont disponibles et, avec des pondérations linéaires, ils portent déjà environ 56~\% du
poids. Reproduire la prime de la plateforme à partir de son carnet d'ordres (notionnel d'impact, échantillonnage toutes les 5 secondes), calculer la moyenne pondérée obtenue jusqu'ici,
prévoir les échantillons restants (niveau actuel, retour à la moyenne) et appliquer le clamp et les plafonds~;
la zone morte rend la réponse exacte dès que la prévision reste à l'intérieur.
\iqlookfor{reproduire l'indice et les pondérations, puis prévoir ce qui reste.}
\end{solution}
